%% notes-tex/wet-lab/040-inhibitor-synergy-wetlab/sections/3-methods.tex

\section{Methods\secstatus{todo}}
\label{sec:methods}

\subsection{Single-agent curves\secstatus{todo}}
\label{sec:hill}

Index the six inhibitors by $i$ and write $d_i$ for the dose of inhibitor $i$ in a well,
in g\,L$^{-1}$. The fitness of a well holding inhibitor $i$ alone is modeled by a Hill
curve with three parameters, a top $\theta_i$, a half-maximal dose $c_i$ and a slope
$h_i$,
%
\begin{equation}
  f_i(d) \;=\; \frac{\theta_i}{1 + (d / c_i)^{h_i}} ,
  \qquad
  \varphi_i(d) \;=\; \frac{f_i(d)}{\theta_i} \;=\; \frac{1}{1 + (d / c_i)^{h_i}} ,
  \label{eq:hill}
\end{equation}
%
where $\varphi_i(d) \in (0, 1)$ is the \term{fractional fitness}, the fitness relative to
the curve's own top. Fitting the top rather than dividing by the controls is what the
ex21 control anomaly of Section~\ref{sec:data} forces. Inverting
Equation~\ref{eq:hill} gives the dose of inhibitor $i$ alone that leaves fractional
fitness $y$,
%
\begin{equation}
  D_i(y) \;=\; c_i \left( \frac{1 - y}{y} \right)^{1 / h_i} ,
  \qquad 0 < y < 1 ,
  \label{eq:inverse}
\end{equation}
%
so $c_i = D_i(1/2)$ is the IC50 and $D_i(0.7)$ is the IC30. Parameters are fit by least squares on the dosed wells of one run, with $h_i$ bounded
above at 20. The primary variant puts a well that did not grow at fitness 0 and the
\texttt{excluded} variant drops it; intervals are the 2.5 and 97.5 percentiles of 1000
resamples of the replicate wells within each dose. Two properties of the fits matter
downstream. A slope at its bound, $h_i = 20$, marks a cliff, fitness falling from near
the top to none between two adjacent tested doses, so $c_i$ is located only within that
dose interval and the bootstrap understates it, since resampling cannot move a cliff that
sits between doses. A fitted $c_i$ at or above the highest dose tested is an
extrapolation of the curve rather than a reading of it.

\subsection{The three reference rules\secstatus{todo}}
\label{sec:rules}

Let $S$ be the set of inhibitors present in a combination, each at its own dose $d_i$,
and let $E$ denote a predicted fractional fitness for that combination. Every rule below
takes the single-agent curves of Section~\ref{sec:hill} and no fitted parameter of its
own.

\notesec{Bliss independence} Bliss 1939 treats the agents as independent events acting
on survival, so the surviving fraction multiplies,
%
\begin{equation}
  E^{\text{Bliss}}_{S} \;=\; \prod_{i \in S} \varphi_i(d_i) ,
  \label{eq:bliss}
\end{equation}
%
which for a pair $A$ and $B$ is the familiar $f_{AB} = f_A f_B$. Two variants differ only
in where $\varphi_i(d_i)$ comes from: \texttt{bliss\_ex21} reads it off the ex21 Hill fit
at the combination's dose, and \texttt{bliss\_ex23} uses the observed ex23 single-agent
wells at that dose, the mean of three wells with a well that did not grow counted as 0.
The second is the stronger reference, since it carries the run's own singles rather than
another run's curve.

\notesec{Loewe additivity} Loewe additivity asks what dose of each agent alone would
produce the combination's effect, and calls the agents additive when those doses sum to
one full dose. With $D_i$ from Equation~\ref{eq:inverse}, the prediction is the effect
$E$ solving
%
\begin{equation}
  \sum_{i \in S} \frac{d_i}{D_i(E)} \;=\; 1 .
  \label{eq:loewe}
\end{equation}
%
The left side is monotone in $E$, because $D_i(E)$ grows without bound as $E \to 0$ and
vanishes as $E \to 1$, so the root is unique and is found numerically on
$\operatorname{logit} E$ by bisection. Evaluated at an observed fractional fitness $y$
instead, the same sum is the \term{combination index}
%
\begin{equation}
  \mathrm{CI}(y) \;=\; \sum_{i \in S} \frac{d_i}{D_i(y)} ,
  \label{eq:ci}
\end{equation}
%
below 1 for synergy and above 1 for antagonism relative to Loewe. A component whose
$D_i(E)$ lies above the highest dose ex21 tested is flagged: that single never reaches
$E$ within its measured range, so its term is a Hill extrapolation.

\notesec{Highest single agent} $E^{\text{HSA}}_{S} = \min_{i \in S} \varphi_i(d_i)$, the
sickest single agent, in the ex21-fit and observed-ex23 variants.

\subsection{Scoring the combinations\secstatus{todo}}
\label{sec:scoring}

\notesec{The growth threshold} A rule predicts fitness, and a well is recorded as grown
or not, so the two are joined by a threshold $\tau$ taken from the data: the lowest
fitness of any grown inhibited ex23 well under that call, which is the slowest growth the
call registers within 85 h. A rule predicts growth when its prediction is at least
$\tau$. The thresholds are 0.248 for the served call and 0.171 for the software call.
Classification is scored as accuracy and as AUROC with the rule's prediction as the
score, so AUROC reads the ordering and is free of $\tau$.

\notesec{Fitness of the grown combinations} Observed fitness is the mean of the grown
wells, the definition experiment 039 used. Rules are scored by Spearman correlation, by
RMSE and by the mean of observed minus predicted, whose sign says which way a rule errs,
over all 63 combinations and over the 57 holding two or more inhibitors; the rows by
number of inhibitors are in \file{results/mixture\_scores.csv}.

\notesec{Pair deviations} For each of the 15 pairs, the deviation is observed fitness
minus $E^{\text{Bliss}}$ from the ex23 singles, with a 95 percent interval from 2000
resamples of the replicate wells of the pair and of both singles (and of the Hill fit
draws for the ex21 variant). A pair is called synergistic when that interval lies below
zero, antagonistic when it lies above, and additive otherwise.

\notesec{Isoboles} Each grid holds two inhibitors over nine doses each, two plate wells
per cell. The Bliss surface is built empirically from the grid's own axes, so the excess
over Bliss is a within-run quantity, averaged over the 81 interior cells with a 95
percent interval from 2000 resamples of the two wells per cell. The Loewe additive front
is the locus of Equation~\ref{eq:loewe} at $E = \tau$ computed from the grid-axis Hill
fits, and the cells that contradict it are counted in both directions.

\subsection{How deletion profiles compose, measured on public data\secstatus{todo}}
\label{sec:het}

Claim 3 cannot be asked of the inhibitors, because no screen of a deletion pool under any
inhibitor combination exists. Hillenmeyer 2008's heterozygous diploid screen (HET) is the
one public source with two compounds in one well: of its 452 environments in the 033
pooled table, 26 carry two compounds, over 5,825 genes. The response of gene $g$ in
environment $e$ is
$y_{ge} = \log_2(\text{mean control intensity} / \text{treatment intensity})$, so a
positive value is a fitness defect; where the 033 flatten keeps several measurements for
a cell, their mean is used. Each two-compound environment is matched to the
single-compound environment of each partner at the same $\log_{10}$ molar dose, else the
nearest dose with the same generation count. A match is reasonable when the dose is
within one twofold step and the generations agree: 46 matches are exact, 2 within
twofold, and 4 are beyond, every one of those four a tacrolimus pair whose dosed singles
sit 402 to 804-fold away. Those four are flagged and reported separately.

Writing $y_1$ and $y_2$ for the two matched single profiles and $y_{12}$ for the pair's,
four fixed rules and one fitted one are scored on the genes measured in all three
environments: the sum $y_1 + y_2$, the mean $(y_1 + y_2)/2$, the elementwise maximum, the
Bliss product on the fitness scale, and a free linear fit $y_{12} = a y_1 + b y_2 + c$.
On this scale the Bliss product is the sum,
%
\begin{equation}
  -\log_2\!\left( 2^{-y_1} \, 2^{-y_2} \right) \;=\; y_1 + y_2 ,
  \label{eq:blisslog}
\end{equation}
%
and the script confirms it: the largest absolute difference between the two is 0 on all
26 pairs, so the Bliss column carries no separate information and is plotted as the sum.
Fixed rules are scored by $R^2 = 1 - \mathrm{SS}_{\text{res}} / \mathrm{SS}_{\text{tot}}$
as they stand, with no refit, so a rule of the wrong magnitude scores negative; Spearman
is scale-free, which is why the sum and the mean share one Spearman value.

A \term{hit} is a gene with robust $|z| > 2$ within its environment, where
$z = (y - \mathrm{median}) / (1.4826\,\mathrm{MAD})$ over that environment's genes; 10.2
percent of all gene-by-environment cells are hits. An \term{emergent} gene is a hit in
the pair and in neither single, a \term{masked} gene a hit in a single and not in the
pair; the emergent fraction divides by the pair's hits, the masked fraction by the union
of the singles' hits. Two nulls of 200 draws each replace the matched singles by random
ones: any two YPD singles without either partner compound, and the stricter null
restricted to singles sharing a control set with the pair, which all 26 pairs do. The
noise floor is the only near-replicate pair of single environments in the release, two
5-fluorouracil doses within 1 percent of each other.

\subsection{Compound-cold profiles, similarity, and GO enrichment\secstatus{todo}}
\label{sec:profiles}

\notesec{The predictor} Experiment 038's nested ridge on FCFP4 counts, the model that
sets the bar on the corrected store. The fingerprint is the 031 recipe: RDKit Morgan
with radius 2 over 2048 bits, feature atom invariants, counts rather than bits, float32,
from the curated identity table's SMILES. Re-featurizing 35 compounds reproduced every
031 stored row exactly
(\file{results/profile\_fingerprint\_check.csv}), which is what makes a prediction for a
compound outside the 32 comparable to one inside. Formic acid (\texttt{OC=O}) and lactic
acid (\texttt{CC(O)C(=O)O}, no stereocenter stated on the stock sheet) are identity gaps
in the private loader and are assigned here.

The regularization strength, rank and clip are chosen by leave-one-compound-out inside
the fitted pool on the centered target. Furfural and 5-HMF are predicted with themselves
held out entirely (31 fitted); the other four are predicted from a fit on all 32. A
prediction is scored by Spearman correlation against the measured profile, raw and after
centering each gene by its mean over the 32 compounds, which removes the component every
compound shares.

\notesec{Similarity} Two measures per compound pair, each raw and centered: the Spearman
correlation over the genes both profiles carry, and the Jaccard overlap of the 100 most
sensitive genes. Claim 2 tests each against the pair deviation of
Section~\ref{sec:scoring} over the 15 pairs with a 10,000-draw permutation $p$ and a
leave-one-pair-out range, under both growth calls and both profile matrices: the
\term{best-available} matrix, which uses a measured profile wherever its replicate
reliability reaches 0.5, and the \term{all-predicted} matrix, in which every column is a
prediction. The 0.5 reliability threshold is a choice, not a sourced value.

\notesec{GO enrichment} Gene-to-term annotations come from the S288C GFF
\texttt{Ontology\_term} fields of the torchcell genome, propagated to \texttt{is\_a}
ancestors through the GO release of 2024-01-17 whose sha256 is recorded in
\file{results/go\_summary.json}. The background is the 3,587 build-002 Vanacloig genes,
all annotated, and the 2,292 terms with 5 to 500 background genes are tested one-sided
hypergeometric with Benjamini-Hochberg control per set at FDR 0.05. Study sets are the
100 most sensitive genes of each measured and predicted profile, the core nomination set,
and each combination's composed top 100. The control set is the 100 most sensitive genes
of the build-002 gene mean over the 32 compounds, which is what a model with no
compound-specific signal outputs, so a term the prediction shares with the control is
evidence of the generic stress component rather than of the compound.
